\documentclass[11pt]{article}
\usepackage[margin=0.9in]{geometry}
\usepackage{amsmath,amssymb,amsthm,mathtools}
\usepackage[colorlinks=true,allcolors=blue]{hyperref}
\newtheorem{theorem}{Theorem}
\newtheorem{lemma}[theorem]{Lemma}
\newtheorem{proposition}[theorem]{Proposition}
\newtheorem{corollary}[theorem]{Corollary}
\newtheorem{remark}[theorem]{Remark}
\newcommand{\NP}{\operatorname{NP}}
\newcommand{\IP}{\operatorname{IP}}
\newcommand{\Ass}{\operatorname{Ass}}
\newcommand{\Min}{\operatorname{Min}}
\newcommand{\R}{\mathbb R}
\newcommand{\N}{\mathbb N}
\title{The Mendez--Pinto--Villarreal conjecture in unmixed height three}
\author{Henry Zweiman}
\date{September 9, 2026}
\begin{document}
\maketitle
\begin{abstract}
Let $I$ be a monomial ideal whose minimal primary components are
irreducible and whose associated primes all have height three.
We prove that equality of all symbolic and ordinary powers forces the
exponent of each variable to agree in every primary component containing
it. Consequently $I$ is a standard variable weighting of its squarefree
radical, which is itself Simis. This proves the unmixed height-three
case of the Mendez--Pinto--Villarreal conjecture, without restrictions
on the exponents or on the supports of minimal generators.
The proof uses a classical complete-intersection criterion in dimension
one to obtain rigidity on four variables. A putative disagreement then
localizes to five variables and four support configurations. Three are
excluded by symbolic-power witnesses or triangle specializations;
the last admits an explicit obstructing vertex of the irreducible
polyhedron for every choice of exponents.
\end{abstract}

\section{Statement and context}
Let $S=k[x_1,\ldots,x_n]$, where $k$ is an arbitrary field. For an ideal
$I\subsetneq S$ its symbolic powers in this paper are
\[
 I^{(m)}=\bigcap_{\mathfrak p\in\Min(S/I)}(I^mS_{\mathfrak p}\cap S).
\]
The ideal is \emph{Simis} if $I^{(m)}=I^m$ for every integer $m\geq1$.
An irreducible monomial ideal is generated by positive powers of
distinct variables. We assume throughout the main theorem that
\begin{equation}\label{eq:decomp}
 I=\bigcap_{P\in\mathcal H}Q_P,
 \qquad Q_P=(x_i^{a_{P,i}}:i\in P),\quad a_{P,i}\geq1,
\end{equation}
where $\mathcal H$ is a family of distinct three-element subsets of
$\{1,\ldots,n\}$. Thus \eqref{eq:decomp} is a minimal primary
decomposition without embedded primes. Conversely, every monomial
ideal with irreducible minimal primary components and all associated
primes of height three has this form. This spells out the meaning of
minimal irreducible decomposition used here: there is exactly one
irreducible primary component for each associated prime.

For positive integers $w_1,\ldots,w_n$ and a monomial ideal $J$, denote
by $J_w$ the ideal generated by the images of its monomials under
$x_i\mapsto x_i^{w_i}$. This is a standard variable weighting.

\begin{theorem}\label{thm:main}
For an ideal as in \eqref{eq:decomp}, the following are equivalent:
\begin{enumerate}
\item $I$ is Simis.
\item For each $i$, the integers $a_{P,i}$ are equal as $P$ ranges
over the members of $\mathcal H$ containing $i$, and $\sqrt I$ is Simis.
\item $I=(\sqrt I)_w$ for a standard variable weighting $w$, and
$\sqrt I$ is Simis.
\end{enumerate}
\end{theorem}

M\'endez, Vaz Pinto, and Villarreal \cite{MPV}, Conjecture 5.7,
propose the implication from the Simis property to standard weighting
for monomial ideals without embedded primes whose irreducible
decomposition is minimal. Their Corollary 5.5 already proves invariance
of the Simis property under standard weighting; the new implication
in Theorem~\ref{thm:main} is agreement of the primary-component
exponents in unmixed height three.

The height-two case follows from Aslam \cite{Aslam}, Theorem 2.3,
and is also formulated as a case of the conjecture by Bijender and
Kumar \cite{BK}, Theorem 4.9. Bordoloi, Das, and Kumar
\cite{BDK2,BDK3} treat ideals constrained by the supports of their
minimal generators; in particular \cite{BDK3}, Theorem 3.10, treats
support-three ideals, in the setup of that paper. Component height
three and generator support three are different conditions.
Our argument makes no assumption on generator support sizes or on
their equality with the supports of the radical's generators.
The general conjecture remains outside the scope of this paper.

The main mechanism is local: two components containing a discrepantly
weighted variable have union of size at most five. Four-variable
localizations give a rigid complete-intersection structure. The
remaining five-variable configurations can then be excluded uniformly
in all positive exponents.

\section{Monomial operations and dimension-one rigidity}
We first record the operations used in the proof. In this section the
sets in $\mathcal H$ can have arbitrary sizes, but are distinct and
incomparable under inclusion. Write $\mathfrak p_P=(x_i:i\in P)$.
The power $Q_P^m$ is $\mathfrak p_P$-primary: membership of a monomial
depends only on its exponents in $P$. Localization at
$\mathfrak p_P$ removes the other components. Consequently
\begin{equation}\label{eq:symbolic}
 I^{(m)}=\bigcap_{P\in\mathcal H}Q_P^m.
\end{equation}
For later use, the exact membership test is
\begin{equation}\label{eq:floor}
 x^b\in Q_P^m
 \quad\Longleftrightarrow\quad
 \sum_{i\in P}\left\lfloor\frac{b_i}{a_{P,i}}\right\rfloor\geq m.
\end{equation}
Indeed, the right side says that $m$ pure-power factors from $Q_P$
can be chosen to divide $x^b$.

\begin{lemma}\label{lem:operations}
Suppose $I$ is Simis and has the decomposition just described.
\begin{enumerate}
\item For a set $T$ of variables, the ideal in $k[x_i:i\in T]$
obtained by retaining exactly the components with $P\subseteq T$
is Simis, if any components remain. Their exponents are unchanged.
\item Let $\pi:S\longrightarrow k[x_i:i\in T]$ set all other
variables equal to zero. Suppose every $\pi(Q_P)$ is proper and nonzero, and their radicals
are distinct and incomparable after deleting identical duplicate
ideals. Then their intersection $\pi(I)$ is Simis.
\end{enumerate}
\end{lemma}
\begin{proof}
For (1), invert the variables outside $T$. A component disappears
precisely when its support is not contained in $T$. Localization
commutes with powers and finite intersections, so the Simis equalities
\eqref{eq:symbolic} pass to the retained components. They descend to
$k[x_i:i\in T]$ by faithful flatness of Laurent polynomial extension.

For (2), if $L$ is monomial, a monomial in the remaining variables
belongs to $\pi(L)$ exactly when it belongs to $L$. Hence $\pi$
commutes with intersections of monomial ideals; as a surjective
homomorphism it also commutes with ideal powers. Apply $\pi$ to
$I^m=\bigcap_PQ_P^m$. Under the stated hypothesis the resulting
components form a minimal irreducible primary decomposition, so
\eqref{eq:symbolic} identifies the right side with $\pi(I)^{(m)}$.
\end{proof}

Part (2) is deliberately conditional. Setting variables equal to zero
is not asserted to preserve the Simis property without checking the
image decomposition. In both applications below that check is explicit.

We use the following classical local criterion in its ideal case:
a generically complete intersection ideal of grade at least two
in a Cohen--Macaulay local ring is a complete intersection if all
its ordinary-power quotients are Cohen--Macaulay. This is the rank-one
specialization of Branco Correia--Zarzuela \cite{BCZ}, Theorem 6.4.
Its mechanism is Burch's inequality, which forces equimultiplicity,
followed by the fact that a generically complete intersection
equimultiple ideal is a complete intersection. The latter criterion
and its proof also appear in \cite{BCZ2}, Theorem 5.3.

\begin{lemma}\label{lem:dimensionone}
Let $h\geq2$. Suppose $K\subset k[z_1,\ldots,z_{h+1}]$ is Simis
and its minimal primary components are irreducible of height $h$.
Then $K$ is a monomial complete intersection, has at most two minimal
primes, and the exponents of any variable agree in all components
containing it.
\end{lemma}
\begin{proof}
Localize at the homogeneous maximal ideal $\mathfrak m$.
For every $m\geq1$, the equality of symbolic and ordinary powers
expresses $K^m$ as an intersection of powers of its primary components.
Those powers are primary with height-$h$ radicals. There are no
embedded associated primes, so the one-dimensional local quotient
by $K^m$ has depth one and is Cohen--Macaulay. At every minimal prime,
$K$ is generated by the $h$ pure powers in its corresponding
irreducible component. It is therefore generically a complete
intersection. The local ring is regular and $K$ has grade $h\geq2$,
so the preceding classical criterion applies.

The number of minimal monomial generators is unchanged on localizing
at $\mathfrak m$, as is seen from $K/\mathfrak mK$. Thus $K$ has
exactly $h$ minimal monomial generators. Their supports are pairwise
disjoint. In fact, if two shared a variable, that variable together
with one variable from each of the other $h-2$ generators would give
a prime of height at most $h-1$ containing $K$, a contradiction.
These $h$ disjoint nonempty supports occupy at most $h+1$ variables.
All therefore have size one, with at most one exception of size two.
The minimal primes select one variable from each generator support,
giving at most two choices.

At such a prime, all unselected variables are units. The localized
ideal is generated by the powers of the selected variables occurring
in their unique monomial generators. Contraction gives the unique
isolated primary component. In particular the exponent of a selected
variable is independent of the prime in which it is selected.
\end{proof}

\begin{corollary}\label{cor:triangle}
An intersection of three irreducible primary ideals with radicals
$(x,y)$, $(x,z)$, and $(y,z)$ in $k[x,y,z]$ is not Simis.
\end{corollary}
\begin{proof}
Lemma~\ref{lem:dimensionone} with $h=2$ permits at most two minimal
primes.
\end{proof}

\section{An obstruction on five variables}
We give a uniform obstruction needed for the last support configuration.
For a monomial ideal $L$ let $\NP(L)$ be its Newton polyhedron,
the convex hull of its exponent set. Equivalently it is the convex
hull of its minimal generator exponents plus $\R_{\geq0}^n$.
For a decomposition into irreducible components, define
\[
 \IP(L)=\bigcap_P\NP(Q_P)
 =\left\{b\in\R_{\geq0}^n:
       \sum_{i\in P}\frac{b_i}{a_{P,i}}\geq1\text{ for every }P\right\}.
\]
These irreducible polyhedra and their relation to integral closures
are studied in Grisalde--Seceleanu--Villarreal \cite{GSV}, Section 7.
Only the following necessary condition is required here.

\begin{lemma}\label{lem:polyhedral}
Under the decomposition assumptions of Lemma~\ref{lem:operations},
if $L$ is Simis, then $\NP(L)=\IP(L)$. Every vertex of $\IP(L)$
is then the exponent vector of a monomial in $L$.
\end{lemma}
\begin{proof}
The inclusion $\NP(L)\subseteq\IP(L)$ follows from $L\subseteq Q_P$.
Take a rational $b\in\IP(L)$ and a positive integer $N$ such that
every $Nb_i/a_{P,i}$ is integral for every relevant pair $(P,i)$,
and every $Nb_i$ is integral. Formula~\eqref{eq:floor} gives
$x^{Nb}\in\bigcap_PQ_P^N=L^N$. Since
$\NP(L^N)=N\NP(L)$, we obtain $b\in\NP(L)$. Rational points are
dense in the rational polyhedron $\IP(L)$, and $\NP(L)$ is closed,
which proves equality.

Finally, every vertex of the convex hull of finitely many generator
exponents plus the nonnegative orthant is one of those exponents.
Indeed a nonzero orthant displacement places the point strictly
between its undisplaced point and twice that displacement, while a
nontrivial convex combination contradicts extremality.
\end{proof}

\begin{proposition}\label{prop:four}
For arbitrary positive integers $a,b,c,\alpha,\beta,\gamma,d,e$,
the ideal in $k[x,y,z,u,v]$ with components
\begin{align*}
 Q_0&=(x^a,y^b,z^c),& Q_1&=(x^\alpha,u^d,v^e),\\
 Q_2&=(y^\beta,u^d,v^e),& Q_3&=(z^\gamma,u^d,v^e)
\end{align*}
is not Simis.
\end{proposition}
\begin{proof}
The four radicals are distinct incomparable triples. Use coordinates
$(X,Y,Z,U,V)$ and put
$r=\alpha/a+\beta/b+\gamma/c$. The irreducible polyhedron has
nonnegative coordinates and inequalities
\begin{equation}\label{eq:fourrows}
 \frac Xa+\frac Yb+\frac Zc\geq1,\quad
 \frac X\alpha+\frac Ud+\frac Ve\geq1,\quad
 \frac Y\beta+\frac Ud+\frac Ve\geq1,\quad
 \frac Z\gamma+\frac Ud+\frac Ve\geq1.
\end{equation}
If $r>1$, take
\[
 p=\left(\frac\alpha r,\frac\beta r,\frac\gamma r,
                  d(1-1/r),0\right).
\]
This is feasible. The four equalities in \eqref{eq:fourrows}
together with $V=0$ uniquely determine $p$: the last three give
$X/\alpha=Y/\beta=Z/\gamma=t$, the first gives $rt=1$, and then
$U=d(1-t)$. Hence $p$ is a vertex. Each of
$X/a,Y/b,Z/c$ is strictly between zero and one, since its numerator
is a positive summand of $r$ and the other two summands are positive.

If $r\leq1$, instead take
\[
 p=\bigl(a(1-\beta/b-\gamma/c),\beta,\gamma,0,0\bigr).
\]
Here $X\geq\alpha$, so the point is feasible. The first, third, and
fourth equalities in \eqref{eq:fourrows}, together with $U=V=0$,
uniquely determine it. It is again a vertex. Since $\alpha/a>0$
and $r\leq1$, we have
$0<X<a$, $0<\beta<b$, and $0<\gamma<c$.

In both cases the vertex has all three first coordinates below the
respective pure-power thresholds of $Q_0$. It is outside the union
of the exponent orthants of $Q_0$, hence cannot be the exponent
vector of a monomial in $\bigcap_{j=0}^3Q_j$.
Lemma~\ref{lem:polyhedral} excludes the Simis property.
\end{proof}

\section{Proof of the height-three theorem}
Assume $I$ is Simis. If two component supports $P,R$ intersect in
two variables, their union has size four. Retain the components
contained in that union by Lemma~\ref{lem:operations}(1), and apply
Lemma~\ref{lem:dimensionone} with $h=3$. Their weights agree on
both common variables. We will use this conclusion repeatedly.

Suppose some variable nevertheless has different exponents in two
components. Those supports must intersect in exactly that variable.
After relabeling, and localizing to their union, write
\[
 P=\{x,a,b\},\qquad R=\{x,c,d\},
\]
with unequal $x$-exponents $A$ and $B$. Here $x,a,b,c,d$ denote
variables; exponent letters in Proposition~\ref{prop:four} will
only be used after a separate relabeling. Every additional support
containing $x$ would have the form $\{x,s,t\}$ with
$s\in\{a,b\}$ and $t\in\{c,d\}$. Its two-variable overlaps with
$P$ and $R$ would force its $x$-exponent to equal both $A$ and $B$.
Thus $P$ and $R$ are the only supports containing $x$.

Call the other supports \emph{extras}. They are triples in
$\{a,b,c,d\}$. There are at most two: if there are any, invert $x$
and apply Lemma~\ref{lem:dimensionone} to the retained height-three
components in four variables. Up to exchanging the pairs
$\{a,b\},\{c,d\}$ and permuting within each pair, there are four
possibilities.

\subsection*{No extras}
Exchange $P,R$ if necessary so that $A<B$, and let $C$ be the
$c$-exponent in $Q_R$. The monomial
\[
 f=x^{\max(2A,B)}c^C
\]
lies in $Q_P^2\cap Q_R^2$. In the subring $k[x,c]$ we have
\[
 I\cap k[x,c]=(x^A)\cap(x^B,c^C)=(x^B,x^Ac^C).
\]
A factor of a monomial involving only $x,c$ also involves only $x,c$,
so ordinary powers commute with this monomial contraction. The three
generator products in degree two have exponent pairs
$(2B,0)$, $(A+B,C)$, and $(2A,2C)$. None divides $f$, because
\[
 \max(2A,B)<2B,\qquad \max(2A,B)<A+B,\qquad C<2C.
\]
Therefore $f\in I^{(2)}\setminus I^2$, a contradiction.

\subsection*{One extra}
By symmetry the extra has support $\{b,c,d\}$. Set $a=d=0$.
The three image components have supports $\{x,b\}$, $\{x,c\}$,
and $\{b,c\}$, respectively. Their radicals are distinct and
incomparable. Lemma~\ref{lem:operations}(2) would make their
intersection Simis, contrary to Corollary~\ref{cor:triangle}.

\subsection*{Two extras omitting variables from different pairs}
The extras can be written $\{b,c,d\}$ and $\{a,b,d\}$.
They share the two variables $b,d$, so the four-variable conclusion
at the start of the proof makes their exponents on $b,d$ equal.
Set $a=c=0$. The image supports are
\[
 \{x,b\},\quad\{x,d\},\quad\{b,d\},\quad\{b,d\}.
\]
The last two ideals are identical, including their exponents.
After deleting one duplicate, Lemma~\ref{lem:operations}(2)
again gives a Simis triangle, contradicting
Corollary~\ref{cor:triangle}.

\subsection*{Two extras omitting variables from the same pair}
The extras can be written $\{a,c,d\}$ and $\{b,c,d\}$.
Together with $R$, they all contain $c,d$. Pairwise four-variable
localization makes the exponents on $c,d$ agree across all three.
After relabeling the variables $(x,a,b,c,d)$ as $(x,y,z,u,v)$,
the four components have exactly the form of
Proposition~\ref{prop:four}. That proposition supplies the final
contradiction.

These possibilities exhaust the subsets of the four available extra
triples of cardinality at most two. Thus every shared-variable
exponent agrees. Let $w_i$ be this common exponent, choosing $w_i=1$
for any variable occurring in no component. Weighting commutes with
monomial intersections, by the least-common-multiple formula for
their generators. Hence
\[
 I=\bigcap_{P\in\mathcal H}(x_i^{w_i}:i\in P)
   =\left(\bigcap_{P\in\mathcal H}(x_i:i\in P)\right)_w
   =(\sqrt I)_w.
\]
Corollary 5.5(b) of \cite{MPV} now makes $\sqrt I$ Simis.
Conversely the same established weighting invariance proves that
every standard weighting of a squarefree Simis ideal is Simis.
Uniqueness of isolated primary components identifies the component
exponents of such a weighting, so the three conditions of
Theorem~\ref{thm:main} are equivalent.\qed

\section{Scope and further questions}
The theorem removes the nonsquarefree weighting question in the
entire unmixed height-three class under the irreducibility hypothesis.
It leaves the separate squarefree Simis classification intact.
It also does not assert that equality of just the second powers
forces the conclusion: the dimension-one criterion uses all powers.

There is a general localization principle behind the argument.
If all component heights are at most $h$, a discrepancy on a shared
variable is retained by localizing to the union of the two relevant
supports, which has at most $2h-1$ variables. Thus a counterexample
to weighting rigidity in bounded height would already occur in
that many variables. In unmixed height three the four-variable
complete-intersection lemma makes the remaining support list short.
For mixed heights at most three, or unmixed height four, additional
configurations remain to be controlled. They are natural next cases;
the present proof does not settle them.

The accompanying Python script checks the finite support reduction
on five labeled variables and the two rational-vertex formulas using
exact arithmetic. These calculations are supplementary: the proof
above applies to all positive exponents without a finite cutoff.
The manuscript and checks were prepared with OpenAI Codex. The
originating system performed the internal audit; no independent
human verification or peer review is represented.

\begin{thebibliography}{9}
\bibitem{Aslam}
A. Aslam, \emph{Symbolic powers of monomial ideals which are generically
complete intersections}, Mathematical Reports \textbf{16(66)} (2014),
no. 1, 69--81.
\href{https://imar.ro/journals/Mathematical_Reports/Pdfs/2014/1/6.pdf}{Publisher PDF}.

\bibitem{BK}
Bijender and A. Kumar, \emph{Waldschmidt constant of monomial ideals
and Simis ideals}, preprint (2025),
\href{https://arxiv.org/abs/2512.22940v1}{arXiv:2512.22940v1}.

\bibitem{BDK2}
P. Bordoloi, K. K. Das, and R. Kumar, \emph{Support-2 monomial ideals
that are Simis}, Journal of Algebraic Combinatorics \textbf{63}
(2026), no. 3, Paper No. 40.
\href{https://arxiv.org/abs/2504.07045v2}{Preprint version: arXiv:2504.07045v2}.

\bibitem{BDK3}
P. Bordoloi, K. K. Das, and R. Kumar, \emph{The Mendez--Pinto--Villarreal
conjecture for some classes of monomial ideals}, preprint (2026),
\href{https://arxiv.org/abs/2607.25683v1}{arXiv:2607.25683v1}.

\bibitem{BCZ}
A. L. Branco Correia and S. Zarzuela, \emph{Some asymptotic properties
of the Rees powers of a module}, preprint (2004),
\href{https://arxiv.org/abs/math/0408351v1}{arXiv:math/0408351v1}.

\bibitem{BCZ2}
A. L. Branco Correia and S. Zarzuela, \emph{On equimultiple modules},
preprint (2007),
\href{https://arxiv.org/abs/0708.2225v1}{arXiv:0708.2225v1}.

\bibitem{GSV}
G. Grisalde, A. Seceleanu, and R. H. Villarreal, \emph{Rees algebras
of filtrations of covering polyhedra and integral closure of powers
of monomial ideals}, Research in the Mathematical Sciences
\textbf{9} (2022), Paper No. 13.
\href{https://doi.org/10.1007/s40687-021-00310-2}{doi:10.1007/s40687-021-00310-2}.

\bibitem{MPV}
F. O. M\'endez, M. Vaz Pinto, and R. H. Villarreal,
\emph{Symbolic powers: Simis and weighted monomial ideals},
Journal of Algebra and Its Applications \textbf{24} (2025),
nos. 13--14, 2541001.
\href{https://arxiv.org/abs/2402.08833v3}{Cited version: arXiv:2402.08833v3}.
\end{thebibliography}
\end{document}
